\chapter{ภาคผนวก ก — คลังสคริปต์ Python สำหรับการวิเคราะห์ข้อมูลไมโครไบโอมและสถิติเกษตรแม่นยำ}
\label{appendix:a_ai_code}

\noindent คลังสคริปต์ภาษา Python ที่ผ่านการทดสอบและใช้งานจริงสำหรับการวิเคราะห์ข้อมูลไมโครไบโอม การคำนวณสถิติเชิงภูมิศาสตร์ การประมวลผลภาพโรคพืชด้วยคอมพิวเตอร์วิทัศน์ และการสร้างแบบจำลองการเรียนรู้ของเครื่องสำหรับเกษตรแม่นยำ\index{Python Scripts@Python Scripts for Agri}\index{Kriging@Kriging (การประมาณค่าเชิงพื้นที่)}\index{Random Forest@Random Forest}\index{Computer Vision@Computer Vision}\index{NDVI@NDVI}

\section{การประมาณค่าความหนาแน่นจุลินทรีย์ดินเชิงพื้นที่ด้วย Ordinary Kriging}
\begin{lstlisting}[language=Python, caption={การสร้างแผนที่ความหนาแน่นจุลินทรีย์ดินด้วย PyKrige}]
import numpy as np
from pykrige.ok import OrdinaryKriging
import matplotlib.pyplot as plt

def interpolate_microbial_density(x_coords, y_coords, cfu_values, grid_res=100):
    """
    ประมาณค่าความหนาแน่นชีวมวลจุลินทรีย์ดินเชิงพื้นที่ด้วยสถิติเชิงภูมิศาสตร์ Ordinary Kriging
    """
    OK = OrdinaryKriging(
        x_coords, y_coords, cfu_values,
        variogram_model='spherical',
        verbose=False, enable_plotting=False
    )
    
    gridx = np.linspace(min(x_coords), max(x_coords), grid_res)
    gridy = np.linspace(min(y_coords), max(y_coords), grid_res)
    z, ss = OK.execute('grid', gridx, gridy)
    
    print(f"[GEO-STAT] Mean Estimated Biomass Density: {np.mean(z):.2f} x 10^6 CFU/g")
    return gridx, gridy, z
\end{lstlisting}

\section{การคำนวณดัชนี NDVI และ NDRE จากข้อมูลภาพถ่ายพหุสเปกตรัมของโดรน}
\begin{lstlisting}[language=Python, caption={การคำนวณดัชนีสุขภาพแปลงเกษตรเพื่อทำนายผลลัพธ์จุลินทรีย์ดิน}]
def compute_vegetation_indices(nir_band, red_band, rededge_band):
    """
    คำนวณดัชนี NDVI และ NDRE เพื่อประเมินคลอโรฟิลล์และความสมบูรณ์ของพืช
    """
    eps = 1e-6
    ndvi = (nir_band - red_band) / (nir_band + red_band + eps)
    ndre = (nir_band - rededge_band) / (nir_band + rededge_band + eps)
    return ndvi, ndre
\end{lstlisting}

\section{การสร้างแบบจำลอง Machine Learning พยากรณ์ประสิทธิภาพจุลินทรีย์ดิน}
\begin{lstlisting}[language=Python, caption={การพยากรณ์ฟอสฟอรัสละลายได้จากแบคทีเรียละลายฟอสเฟตด้วย Scikit-Learn}]
import numpy as np
import pandas as pd
from sklearn.ensemble import RandomForestRegressor
from sklearn.model_selection import train_test_split
from sklearn.metrics import mean_squared_error, r2_score

def train_pgpr_predictor(sensor_csv_path):
    """
    ฝึกฝนโมเดล Random Forest ทำนายอัตราการปลดปล่อย P จากเซนเซอร์ดิน (pH, EC, Moisture, Temp)
    """
    df = pd.read_csv(sensor_csv_path)
    X = df[['soil_ph', 'soil_ec', 'soil_moisture', 'soil_temp', 'organic_matter']]
    y = df['available_p_ppm']
    
    X_train, X_test, y_train, y_test = train_test_split(X, y, test_size=0.2, random_state=42)
    
    rf = RandomForestRegressor(n_estimators=100, max_depth=8, random_state=42)
    rf.fit(X_train, y_train)
    
    y_pred = rf.predict(X_test)
    r2 = r2_score(y_test, y_pred)
    rmse = np.sqrt(mean_squared_error(y_test, y_pred))
    
    print(f"[ML-EVAL] Model R2 Score: {r2:.4f}, RMSE: {rmse:.2f} ppm")
    return rf
\end{lstlisting}

\section{การตัดแยกพื้นที่แผลโรคพืชบนใบด้วย OpenCV และ HSV Thresholding}
\begin{lstlisting}[language=Python, caption={การคำนวณดัชนี Disease Severity Index (DSI) ด้วย OpenCV}]
import cv2
import numpy as np

def calculate_disease_severity(image_path):
    """
    วิเคราะห์ภาพใบพืช แยกพื้นที่รอยแผลสะเก็ดและคำนวณค่า Disease Severity Index
    """
    img = cv2.imread(image_path)
    hsv = cv2.cvtColor(img, cv2.COLOR_BGR2HSV)
    
    # ช่วงสีของใบไม้เขียวสุขภาพดี
    lower_green = np.array([25, 40, 40])
    upper_green = np.array([85, 255, 255])
    mask_green = cv2.inRange(hsv, lower_green, upper_green)
    
    # ช่วงสีของรอยแผลไหม้/เน่าสีน้ำตาล
    lower_brown = np.array([10, 50, 20])
    upper_brown = np.array([24, 255, 200])
    mask_disease = cv2.inRange(hsv, lower_brown, upper_brown)
    
    green_pixels = cv2.countNonZero(mask_green)
    disease_pixels = cv2.countNonZero(mask_disease)
    total_leaf_pixels = green_pixels + disease_pixels
    
    if total_leaf_pixels == 0:
        return 0.0
        
    dsi = (disease_pixels / total_leaf_pixels) * 100.0
    print(f"[CV-AI] Total Leaf Area: {total_leaf_pixels} px, Lesion: {disease_pixels} px")
    print(f"[CV-AI] Calculated DSI: {dsi:.2f}%")
    return dsi
\end{lstlisting}
